\documentclass{article}
\usepackage{graphicx}
\usepackage{amsmath}
\usepackage{amsfonts}
\usepackage{bookmark}
\usepackage{booktabs}
\usepackage{exercises}
\graphicspath{ {} }
\usepackage{indentfirst}
\usepackage{amsthm}
\usepackage{yhmath}
\usepackage{geometry}
\usepackage{placeins}
\usepackage{csquotes}
\usepackage{pgfplots}
\newcommand*{\oldneg}{\mathord{\sim}}
\newcommand\showdiv[1]{\overline{\smash{)}#1}}
\newcommand\Mydiv[1]{%
\smash{\raise.3ex\hbox{$\big)$}}\mkern-8mu
        \overline{\enspace\strut#1}}
\setlength{\tabcolsep}{5pt}
\theoremstyle{definition}
\newtheorem{definition}{Definition}

\setlength{\parindent}{.5in}
\setlength{\parskip}{12pt}
\renewcommand{\arraystretch}{2}
\geometry{margin=1in}
\newcommand{\?}{\stackrel{?}{=}}
\pgfplotsset{width=7cm,compat=1.18}\usepgfplotslibrary{statistics}

\title{8.1 Answers}
\begin{document}
\maketitle

\section*{1}
The distance formula states, for any two points $(x_1,y_1)$ and $(x_2,y_2)$, the distance between them is \[ \sqrt{(x_2-x_1)^2+(y_2-y_1)^2} \]
The midpoint formula states that the midpoint of the line segment formed between any two points $(x_1,y_1)$ and $(x_2,y_2)$ is
$(\frac{x_1+x_2}{2},\frac{y_1+y_2}{2})$.

\section*{2}
It does not matter, as, regardless of which point you choose, $x_2-x_1=-(x_1-x_2)$ and $y_2-y_1=-(y_1-y_2)$, and thus
\[ \sqrt{(x_2-x_1)^2+(y_2-y_1)^2}=\sqrt{(-(x_1-x_2))^2+(-(y_1-y_2))^2}=\sqrt{(x_1-x_2)^2+(y_1-y_2)^2} \] meaning the points 
$(x_1,y_1)$ and $(x_2,y_2)$ can be switched.

\section*{13}
The distance is \[ \sqrt{(15+5)^2+(8-1)^2}=\sqrt{449} \]
and the midpoint is \[ (\frac{-5+15}{2},\frac{1+8}{2})=(5,4.5) \]

\section*{14}
The distance is \[ \sqrt{(-16-9)^2+(-16-9)^2}=\sqrt{1250}=25\sqrt{2} \]
and the midpoint is \[ (\frac{9-16}{2},\frac{9-16}{2})=(-\frac{7}{2},-\frac{7}{2}) \]

\section*{18}
The distance is \[ \sqrt{(6+4)^2+(6-3)^2}=\sqrt{109} \]
C

\section*{19}
The midpoint is \[ (\frac{-3+5}{2},\frac{7-2}{2})=(1,\frac{5}{2}) \]
A

\section*{30}
We are given the three points $(3,5),(6,9),(11,9)$ which form a triangle. Let $A=(3,5),B=(6,9),C=(11,9)$. Then:
\[ AB=\sqrt{(6-3)^2+(9-5)^2}=\sqrt{25}=5 \]
\[ AC=\sqrt{(11-3)^2+(9-5)^2}=\sqrt{80}=2\sqrt{10} \]
\[ BC=\sqrt{(11-6)^2+(9-9)^2}=\sqrt{25}=5 \]
Thus, because exactly one pair of sides ($AB$ and $BC$) are equal, the triangle is isosceles.

\section*{36}
We are given the points $(5,10),(10,7)$, and are finding the perpendicular bisector of the line segment joining these two points. 
We know that the midpoint of the two points is \[ (\frac{5+10}{2},\frac{10+7}{2})=(\frac{15}{2},\frac{17}{2}) \]
Also, the slope of the lines is $\frac{10-7}{5-10}=-\frac{3}{5}$ meaning the slope of the perpendicular bisector will be 
$\frac{5}{3}$. Thus, we can solve the following equation:
\[ \frac{17}{2}=\frac{5}{3}(\frac{15}{2})+b \]
\[ b=-4 \]
Hence the equation of the perpendicular bisector is $y=\frac{5}{3}x-4$.

\section*{40}
We are given the points $(2,6),(3,1),(7,5)$. The first point, $(2,6)$, is opposite to the line segement from $(3,1)$ to $(7,5)$. 
The midpoint of that segment is $(\frac{3+7}{2},\frac{1+5}{2})=(5,3)$. Hence, we know that the slope of the median is 
$\frac{6-3}{2-5}=-1$. Then:
\[ y-3=-(x-5) \]
\[ y=-x+8 \] which is the equation of the median.

\section*{44}
We know that:
\[ \sqrt{74}=\sqrt{(8-1)^2+(13-y)^2} \]
\[ 74=49+(13-y)^2 \]
\[ 13-y=\pm\sqrt{25} \]
\[ -y=-8\vee -y=-18 \] \[ y=8\vee y=18 \]

\section*{45}
Suppose $(x,y)$ is any point on the line $y=2x$. Thus $(x,y)=(x,2x)$. Hence the distance between $(x,y)$ and $(2,3)$ is
\[ D=\sqrt{(2-x)^2+(3-2x)^2} \]
We can solve the above equation when $D=\sqrt{10}$ to find the answer to the last question:
\[ \sqrt{10}=\sqrt{(2-x)^2+(3-2x)^2} \]
\[ 10=(2-x)^2+(3-2x)^2 \]
\[ (x^2-4x+4)+(4x^2-12x+9)=10 \]
\[ 5x^2-16x+13=10 \]
\[ 5x^2-16x+3=0 \]
\[ (5x-1)(x-3)=0 \]
\[ x=\frac{1}{5}\vee x=3 \]
Meaning two points on the line $y=2x$ that are each $\sqrt{10}$ units from $(2,3)$ are 
$(\frac{1}{5},\frac{2}{5})$ and $(3,6)$.

\section*{46}
Consider the point $M(\frac{x_1+x_2}{2},\frac{y_1+y_2}{2})$.
The distance from this to $(x_1,y_1)$ is:
\[ \sqrt{(\frac{x_1+x_2}{2}-x_1)^2+(\frac{y_1+y_2}{2}-y_1)^2} \] 
\[ =\sqrt{(\frac{x_2-x_1}{2})^2+(\frac{y_2-y_1}{2})^2} \] 
\[ =\sqrt{\frac{x_2^2+x_1^2-2x_1x_2}{4}+\frac{y_2^2+y_1^2-2y_1y_2}{4}} \]
and the distance from this point to $(x_2,y_2)$ is:
\[ \sqrt{(\frac{x_1+x_2}{2}-x_2)^2+(\frac{y_1+y_2}{2}-y_2)^2} \]
\[ =\sqrt{(\frac{x_1-x_2}{2})^2+(\frac{y_1-y_2}{2})^2} \]
\[ =\sqrt{\frac{x_2^2+x_1^2-2x_1x_2}{4}+\frac{y_2^2+y_1^2-2y_1y_2}{4}} \]
meaning $M$ is equidistant from $(x_1,y_1)$ and $(x_2,y_2)$. Now we need to show that $M$ really lies on 
the line segment between $(x_1,y_1)$ and $(x_2,y_2)$. The slope between these two points is
\[ \frac{y_1-y_2}{x_1-x_2} \] and the line connecting them is:
\[ y-y_1=\frac{y_1-y_2}{x_1-x_2}(x-x_1) \]
\[ y=\frac{y_1-y_2}{x_1-x_2}(x-x_1)+y_1 \]
and thus
\[ y=\frac{y_1-y_2}{x_1-x_2}(\frac{x_1+x_2}{2}-x_1)+y_1 \]
\[ =\frac{y_1-y_2}{x_1-x_2}\frac{x_2-x_1}{2}+y_1 \]
\[ =\frac{(y_1-y_2)(x_2-x_1)}{2(x_1-x_2)}+y_1 \]
\[ =\frac{-(y_1-y_2)}{2}+y_1 \]
\[ =\frac{y_1+y_2}{2} \] thus confirming $M$ really is on the line segment connecting the two points.

\section*{55}
We are given three points on the circle: $(-220,220),(0,0)$ and $(200,40)$. To find the center of the circle, we need to find the perpendicular bisectors of the segments connecting 
$(-220,220)$ with $(0,0)$ and $(0,0)$ and $(200,40)$ and then find where they intersect. The midpoint of the segments are 
$(\frac{-220+0}{2},\frac{220+0}{2})=(-110,110)$ and $(\frac{0+200}{2},\frac{0+40}{2})=(100,20)$, respectively. They have slopes of 
$\frac{0-220}{0+220}=-1$ and $\frac{40-0}{200-0}=\frac{1}{5}$, respectively. Thus the slopes of the their perpendicular bisectors will be 
$1$ and $-5$, respectively. Hence the equations for their perpendicular bisectors will be:
\[ y-110=x+110 \] \[ y=x+220 \] and \[ y-20=-5(x-100) \] \[ y=-5x+520 \] respectively. We now find where these two lines intersect:
\[ y=-5x+520 \] \[ x+220=-5x+520 \] \[ 6x=300 \] \[ x=50 \] and $y=50+220=270$. Thus the center of the circle is at 
$(50,270)$. This means the radius of the circle is 
\[ \sqrt{(50-0)^2+(270-0)^2}=\sqrt{75400} \] meaning the diameter is about 
$2\sqrt{75400}\approx 550$ feet.


\section*{56}
(a): \[ M_1=(\frac{x_1+0}{2},\frac{y_1+0}{2})=(\frac{x_1}{2},\frac{y_1}{2}) \] and
\[ M_2=(\frac{x_1+x_2}{2},\frac{y_1+0}{2})=(\frac{x_1+x_2}{2},\frac{y_1}{2}) \]


(b): The lenth of the strip is:
\[ \sqrt{(\frac{x_1+x_2}{2}-\frac{x_1}{2})^2+(\frac{y_1}{2}-\frac{y_1}{2})^2} \]
\[ =\sqrt{(\frac{x_2}{2})^2} \]
\[ =\frac{x_2}{2} \]
The sail's base has a length of $x_2$, meaning the length of the strip is half of the base. 


(c): These results do not depend upon the shape of the sail. This can be shown by generalzing the results obtained in parts (a) and (b). Consider any triangle 
$A(x_1,y_1),B(x_2,y_2),C(x_3,y_3)$. Suppose we took the base to be $AC$. This has a distance of:
\[ \sqrt{(x_3-x_1)^2+(y_3-y_1)^2} \]
The midpoint $M_1$ of $AB$ is: \[ (\frac{x_1+x_2}{2},\frac{y_1+y_2}{2}) \] and the midpoint $M_2$ of $BC$ is \[ (\frac{x_2+x_3}{2},\frac{y_2+y_3}{2}) \]
The distance between $M_1$ and $M_2$ is:
\[ \sqrt{(\frac{x_2+x_3}{2}-\frac{x_1+x_2}{2})^2+(\frac{y_2+y_3}{2}-\frac{y_1+y_2}{2})^2} \]
\[ =\sqrt{(\frac{x_3-x_1}{2})^2+(\frac{y_3-y_1}{2})^2} \]
\[ =\sqrt{\frac{(x_3-x_1)^2}{4}+\frac{(y_3-y_1)^2}{4}} \]
\[ =\sqrt{\frac{(x_3-x_1)^2+(y_3-y_1)^2}{4}} \]
\[ =\dfrac{\sqrt{(x_3-x_1)^2+(y_3-y_1)^2}}{2} \]
Thus confirming the distance between $M_1$ and $M_2$ really is half of the length of the base.

\section*{57}
\begin{samepage}
We know that the car is traveling (at time $t=0$) 60 miles an hour from a point $100$ miles west and $40$ miles north from a radio tower, and that the radio tower has a range of 50 miles. Thus, at time $t=0$, 
the distance between the car and the radio tower is \[ \sqrt{(100-0)^2+(-40-0)^2}=\sqrt{11600}\approx 107.7 \] and at time $t=1$, the distance between the car and the radio tower is 
\[ \sqrt{(100-60)^2+(-40-0)^2}\approx 56.6 \] and, in general, at time $t$, the distance between the car and the radio tower is (because the car never moves north or south):
\[ \sqrt{(100-60t)^2+(-40-0)^2} \] Thus we solve the following equation:
\[ \sqrt{(100-60t)^2+(-40)^2}\le 50 \]
\[ (100-60t)^2+1600\le 2500 \]
\[ 3600t^2-12000t+11600\le 2500 \]
\[ 400(9t^2-30t+29)\le 2500 \]
\[ 9(t^2-\frac{10}{3}t)\le \frac{25}{4}-29 \]
\[ 9(t^2-\frac{10}{3}t+(\frac{5}{3})^2)\le \frac{25}{4}-29+9(\frac{5}{3})^2 \]
\[ 9(t-\frac{5}{3})^2\le \frac{9}{4} \]
\[ t-\frac{5}{3}\le\sqrt{\frac{1}{4}}\wedge t-\frac{5}{3}\ge -\sqrt{\frac{1}{4}} \]
\[ t\le \frac{5}{3}+\frac{1}{2}\wedge t\ge \frac{5}{3}-\frac{1}{2} \]
\[ t\le \frac{13}{6}\wedge t\ge\frac{7}{6} \]
Thus, the times $t$ during which the car is in range of the tower are $\frac{7}{6}\le t\le \frac{13}{6}$.
\end{samepage}

\end{document}
